\section{Why This Chapter Must Avoid Fake Rigor}

The manuscript needs a formal language because later chapters discuss state
formation, state transition, and intervention design. But the language must be
honest about what kind of formalism it is. At present, the project does
\emph{not} have a fitted quantitative model with experimentally identified
coefficients for human daily behaviour. It therefore should not present
arbitrary formulas as if they were already established mathematics or settled
empirical law.

This chapter takes a stricter position. It provides:
\begin{enumerate}[nosep]
  \item a disciplined vocabulary for describing behavioural states;
  \item a heuristic scoring scaffold for comparing sources of resistance;
  \item qualitative monotonic assumptions about barriers and preparation;
  \item explicit evidence boundaries for what this notation does and does not
        claim.
\end{enumerate}

\begin{discussion}
The source support for this chapter comes from two directions. The
\textbf{textbook foundation} comes from dynamical-systems texts such as Eugene
Izhikevich's \emph{Dynamical Systems in Neuroscience}, Michael Brin and Garrett
Stuck's \emph{Introduction to Dynamical Systems}, and Steven Strogatz's
\emph{Nonlinear Dynamics and Chaos}. These sources justify using state-space,
stability, and transition language as mathematical scaffolding. The
\textbf{biological legitimacy} comes from later empirical chapters and their
sources, which show that structured state transitions and stability are not
mere metaphors. What they do \emph{not} supply is a calibrated scalar law for
everyday action failure.
\end{discussion}
